\documentclass[10pt]{article}
\usepackage[a4paper,margin=0.65in]{geometry}
\usepackage{amsmath,amssymb,amsfonts,mathtools}
\usepackage[dvipsnames]{xcolor}
\usepackage{array,booktabs,float}
\usepackage[colorlinks=true,linkcolor=blue]{hyperref}
\allowdisplaybreaks
\numberwithin{equation}{section}
\setlength{\jot}{4pt}
\setlength{\emergencystretch}{3em}
\newcommand{\KK}{\mathcal K}
\newcommand{\Dh}{\hat\Delta}
\newcommand{\OO}{\mathcal O}
\newcommand{\LO}{\frac{d^3N_{\rm LO}}{d^3k}}
\newcommand{\new}[1]{{\color{blue}#1}}
\newcommand{\old}[1]{{\color{red}#1}}

\title{Group II, Row IV: fragmentation function and DGLAP}
\date{}

\begin{document}
\maketitle

%==========================================================================================
\section{Your equation}\label{sec:check}
Your $\xi$ form is correct. One change is needed, in blue: the measure is $\int_{x,x',y,y'}$ in both terms, not $\int_{x,y,y'}$. The charge $x'$ of the conjugate amplitude is integrated, as at LO.
\begin{align}
\frac{d^{3}N_{\rm NLO}}{d^{3}k}={}&\frac{1}{(2\pi)^{3}}\,\frac{g^{4}d_{\perp}}{4\pi^{3}}\,\frac{N_{c}}{k^{+}}\int_{\frac{k^{+}}{k^{+}+\Lambda}}^{\frac{k^{+}}{\vee}}\frac{d\xi}{2\pi}\int_{x,\new{x'},y,y'}e^{-ik\cdot(y'-y)\frac{1}{\xi}}\,\frac{(x-y)\cdot(x'-y')}{(x-y)^{2}(x'-y')^{2}}\Big[\xi(1-\xi)+\frac{1-\xi}{\xi}+\frac{\xi}{1-\xi}\Big]\frac{1}{\xi^{2}}\notag\\
&\times\Big[\frac{1}{\epsilon}+\gamma_{E}+\frac{1}{2}\log\frac{(y'-y)^{2}\mu^{2}_{\overline{MS}}}{4}\Big]\big[U^{e'c}(y')-U^{e'c}(x')\big]\big[U^{ec}(y)-U^{ce}(x)\big]\rho^{e'}(x')\rho^{e}(x)\notag\\
&+\frac{1}{(2\pi)^{3}}\,\frac{g^{4}}{4\pi^{3}}\,\frac{N_{c}}{k^{+}}\int_{\frac{k^{+}}{k^{+}+\Lambda}}^{\frac{k^{+}}{\vee}}\frac{d\xi}{2\pi}\int_{x,\new{x'},y,y'}e^{-ik\cdot(y'-y)\frac{1}{\xi}}\,\frac{1}{\xi^{2}}\Bigg[\frac{(x-y)\cdot(y'-y)\,(x'-y')\cdot(y'-y)}{(x-y)^{2}(x'-y')^{2}(y'-y)^{2}}\,d_{\perp}\,\xi(1-\xi)\notag\\
&+\frac{(x-y)\cdot(x'-y')}{(x-y)^{2}(x'-y')^{2}}\Big[\frac{1-\xi}{\xi}+\frac{\xi}{1-\xi}\Big]\Bigg]\big[U^{e'c}(y')-U^{e'c}(x')\big]\big[U^{ec}(y)-U^{ce}(x)\big]\rho^{e'}(x')\rho^{e}(x).
\label{eq:user}
\end{align}
It is Eq.~(1.5) of \texttt{RowIV\_dperp.pdf}, term by term (\texttt{sympy}: difference $0$):
\begin{itemize}
\item the limits are reversed, $\int_{k^+/(k^++\Lambda)}^{k^+/\vee}=-\int_{k^+/\vee}^{k^+/(k^++\Lambda)}$, which removes the overall minus sign of the $E$ form of (1.5);
\item $\frac{1}{\epsilon}+\gamma_E+\frac12\log\frac{(y'-y)^2\mu^2}{4}=\frac{E}{2}$, and $\frac{g^4d_\perp}{4\pi^3}\cdot\frac{1}{2\pi}\cdot\frac{E}{2}=\frac{g^4}{16\pi^4}\,d_\perp E$;
\item in the second term, $\frac{g^4}{4\pi^3}\cdot\frac{1}{2\pi}=\frac{g^4}{16\pi^4}\cdot2$, matching the $2\big(\frac{\xi}{1-\xi}+\frac{1-\xi}{\xi}\big)$ and $2d_\perp\xi(1-\xi)$ of (1.5);
\item the $\frac{1}{\xi^2}$ and the phase $e^{-ik\cdot(y'-y)/\xi}$ are right.
\end{itemize}

\paragraph{Notation used below.}
\begin{itemize}
\item $\Delta\equiv y'-y$, $\Dh\equiv\Delta/|\Delta|$, $X\equiv\frac{x-y}{(x-y)^2}$, $X'\equiv\frac{x'-y'}{(x'-y')^2}$, and
\begin{equation}
\OO\equiv\big[U^{ce'}(y')-U^{ce'}(x')\big]\rho^{e'}(x')\,\big[U^{ce}(y)-U^{ce}(x)\big]\rho^{e}(x),
\end{equation}
which is your colour factor, since $U^{bc}=U^{cb}$.
\item $\xi_0\equiv\frac{k^+}{\vee}$ and $\xi_1\equiv\frac{k^+}{k^++\Lambda}$. The light-cone variables are $\delta Y=\log\frac{\vee}{\Lambda}$ and $\delta\equiv\frac{\Lambda}{k^+}$.
\item The pole is written as $\frac1{\bar\epsilon}\equiv-\frac2\epsilon$, i.e.\ $\frac{1}{\bar\epsilon}=-\frac{1}{\epsilon_{\rm std}}$ in the usual $D=4-2\epsilon_{\rm std}$ notation. Then $\frac{E}{2}=-\frac{\bar E}{2}$, with
\begin{equation}
\bar E\equiv\frac1{\bar\epsilon}-\ell,\qquad \ell\equiv\log\frac{\Delta^2\mu^2_{\overline{MS}}e^{2\gamma_E}}{4}=\log\frac{\Delta^2\mu^2_{\overline{MS}}}{c_0^2},\qquad c_0\equiv2e^{-\gamma_E}.
\end{equation}
The pole is collinear (infrared): $\epsilon<0$, so $\frac1{\bar\epsilon}>0$.
\item The real part of the gluon splitting function is
\begin{equation}
P^{(r)}_{gg}(\xi)\equiv2N_c\Big[\xi(1-\xi)+\frac{\xi}{1-\xi}+\frac{1-\xi}{\xi}\Big]=2N_c\hat\Phi(\xi).
\end{equation}
\end{itemize}

%==========================================================================================
\section{Row IV as a convolution with LO}\label{sec:conv}
\paragraph{LO and two weighted copies of it.} The LO spectrum is
\begin{equation}
\LO(k^+,\mathbf k)=\frac{1}{(2\pi)^3}\frac{g^2}{2\pi^2k^+}\int_{x,x',y,y'}e^{-ik\cdot(y'-y)}\,\frac{(x-y)\cdot(x'-y')}{(x-y)^2(x'-y')^2}\,\OO .
\label{eq:LO}
\end{equation}
Insert a weight into the integrand of Eq.~\eqref{eq:LO}, with the same prefactor:
\begin{align}
\mathcal L(k^+,\mathbf k)&\equiv\frac{1}{(2\pi)^3}\frac{g^2}{2\pi^2k^+}\int_{x,x',y,y'}e^{-ik\cdot(y'-y)}\,\frac{(x-y)\cdot(x'-y')}{(x-y)^2(x'-y')^2}\,\frac{d_\perp\bar E-2}{2}\,\OO ,\label{eq:calL}\\
\mathcal L_D(k^+,\mathbf k)&\equiv\frac{1}{(2\pi)^3}\frac{g^2}{2\pi^2k^+}\int_{x,x',y,y'}e^{-ik\cdot(y'-y)}\Big[\frac{(x-y)\cdot(x'-y')}{(x-y)^2(x'-y')^2}-d_\perp\frac{(x-y)\cdot(y'-y)\,(x'-y')\cdot(y'-y)}{(x-y)^2(x'-y')^2(y'-y)^2}\Big]\OO .\label{eq:calLD}
\end{align}
$\bar E$ depends on $|y'-y|$, so $\mathcal L$ is not a multiple of LO. Its pole part is: $\mathcal L\big|_{\rm pole}=\frac{d_\perp}{2}\frac{1}{\bar\epsilon}\LO$.

\paragraph{The conversion.} At the argument $(\frac{k^+}{\xi},\frac{\mathbf k}{\xi})$, Eq.~\eqref{eq:LO} has prefactor $\frac{g^2\xi}{2\pi^2k^+}$ and phase $e^{-ik\cdot\Delta/\xi}$. Hence, for any weight $W$,
\begin{equation}
\int_{x,x',y,y'}e^{-ik\cdot\Delta/\xi}\,X\cdot X'\,W\,\OO=(2\pi)^3\,\frac{2\pi^2k^+}{g^2\xi}\times\Big[\text{Eq.~\eqref{eq:LO} with weight }W\Big]\Big(\frac{k^+}{\xi},\frac{\mathbf k}{\xi}\Big),
\end{equation}
and the prefactor of Row IV becomes
\begin{equation}
\frac{1}{(2\pi)^3}\frac{g^4N_c}{16\pi^4}\frac{1}{k^+}\frac{1}{\xi^2}\times(2\pi)^3\frac{2\pi^2k^+}{g^2\xi}=\frac{\alpha_s}{2\pi}\,\frac{N_c}{\xi^3}\qquad(\alpha_s=\tfrac{g^2}{4\pi};\ \texttt{sympy}: $0$).
\end{equation}

\paragraph{Result.} With $N_c\,d_\perp\hat\Phi\,\bar E-2N_c\big(\frac{\xi}{1-\xi}+\frac{1-\xi}{\xi}\big)=P^{(r)}_{gg}\frac{d_\perp\bar E-2}{2}+2N_c\xi(1-\xi)$, Row IV is, exactly and for any $d_\perp$,
\begin{equation}
\boxed{\frac{d^3N_{\rm IV}}{d^3k}=\frac{\alpha_s}{2\pi}\int_{\xi_0}^{\xi_1}\frac{d\xi}{\xi^3}\Big\{P^{(r)}_{gg}(\xi)\,\mathcal L\Big(\frac{k^+}{\xi},\frac{\mathbf k}{\xi}\Big)+2N_c\,\xi(1-\xi)\,\mathcal L_D\Big(\frac{k^+}{\xi},\frac{\mathbf k}{\xi}\Big)\Big\}}
\label{eq:conv}
\end{equation}
(\texttt{sympy}: equal to the compact form of \texttt{RowIV\_compact.pdf}, difference $0$).
\begin{itemize}
\item \textbf{The pole} is in $\mathcal L$ only:
\begin{equation}
\frac{d^3N_{\rm IV}}{d^3k}\bigg|_{\rm pole}=\frac{\alpha_s}{2\pi}\,\frac{d_\perp}{2}\,\frac{1}{\bar\epsilon}\int_{\xi_0}^{\xi_1}\frac{d\xi}{\xi^3}\,P^{(r)}_{gg}(\xi)\,\LO\Big(\frac{k^+}{\xi},\frac{\mathbf k}{\xi}\Big).
\label{eq:pole}
\end{equation}
\item \textbf{$\mathcal L_D$ has no pole.} Its tensor $\delta^{ij}-d_\perp\Dh^i\Dh^j$ is traceless in $d_\perp$ dimensions, so it vanishes when averaged over the direction of $\Delta$. It is the linearly polarized part of the $g\to gg$ splitting; it is finite and stays in the coefficient function.
\item \textbf{The measure $d\xi/\xi^3$} is the Jacobian of $\frac{d^3N}{d^3k}=\frac{dN}{dk^+d^2\mathbf k}$: the parent has $(k^+,\mathbf k)/\xi$, so $d^3k_{\rm parent}=d^3k/\xi^3$. For $\frac{dN}{d\eta\,d^2\mathbf k}=k^+\frac{d^3N}{d^3k}$ the same equation has $d\xi/\xi^2$.
\item \textbf{The lower limit $\xi_0=k^+/\vee$} is the condition that the parent momentum $k^+/\xi$ stays below $\vee$.
\end{itemize}

\paragraph{The two choices of $d_\perp$.}
\begin{itemize}
\item $d_\perp=2-\epsilon$ (as in Group I, Rows II and V). Then $\frac{d_\perp\bar E-2}{2}=\frac{1}{\bar\epsilon}-\ell+O(\epsilon)$ (\texttt{sympy}: $0$). The ``$+1$'' from $\frac{d_\perp}{2}\frac1{\bar\epsilon}=\frac1{\bar\epsilon}+1$ cancels the $-1$. So
\begin{equation}
\mathcal L=\frac{1}{\bar\epsilon}\LO-\frac{d^3N_{\ell}}{d^3k},\qquad \mathcal L_D=\LO-2\,\frac{d^3N_D}{d^3k},
\label{eq:CDR}
\end{equation}
where $\frac{d^3N_\ell}{d^3k}$ is Eq.~\eqref{eq:LO} with weight $\ell$ and $\frac{d^3N_D}{d^3k}$ is Eq.~\eqref{eq:LO} with $X\cdot X'\to(X\cdot\Dh)(X'\cdot\Dh)$. The residue of the pole is $P^{(r)}_{gg}$ exactly.
\item $d_\perp=2$. Then $\frac{d_\perp\bar E-2}{2}=\frac1{\bar\epsilon}-\ell-1$. The residue is again $P^{(r)}_{gg}$. The finite part differs by $-\frac{\alpha_s}{2\pi}\int\frac{d\xi}{\xi^3}P^{(r)}_{gg}\,\frac{d^3N_{\rm LO}}{d^3k}$, which is a change of factorization scheme.
\end{itemize}
The DGLAP kernel is the residue, so it is the same in both choices. The finite parts are not, so every row has to be evaluated with the same $d_\perp$.

%==========================================================================================
\section{The endpoint \texorpdfstring{$\xi\to1$}{xi to 1} and the plus prescription (exact limits)}\label{sec:plus}
The upper limit $\xi_1=\frac{k^+}{k^++\Lambda}<1$ regulates the soft gluon, $p^+\to\Lambda$. To compare with DGLAP, use the standard plus distribution on $[0,1]$, cut at $\xi_0$:
\begin{equation}
\int_{\xi_0}^1d\xi\,\Big[\frac{1}{1-\xi}\Big]_+f(\xi)\equiv\int_{\xi_0}^1d\xi\,\frac{f(\xi)-f(1)}{1-\xi}+f(1)\log(1-\xi_0),
\end{equation}
and $P^{(r)}_{gg,+}(\xi)\equiv2N_c\big[\frac{\xi}{(1-\xi)_+}+\frac{1-\xi}{\xi}+\xi(1-\xi)\big]$, with $\frac{\xi}{(1-\xi)_+}\equiv\xi\big[\frac{1}{1-\xi}\big]_+$. Let $f(\xi)$ be any function regular at $\xi=1$. Exactly,
\begin{equation}
\int_{\xi_0}^{\xi_1}d\xi\,P^{(r)}_{gg}(\xi)f(\xi)=\int_{\xi_0}^1d\xi\,P^{(r)}_{gg,+}(\xi)f(\xi)+2N_c\log\frac{1}{1-\xi_1}\,f(1)-\int_{\xi_1}^1d\xi\Big[P^{(r)}_{gg}(\xi)f(\xi)-\frac{2N_cf(1)}{1-\xi}\Big]
\label{eq:plusid}
\end{equation}
(numerically, rel.\ $1\times10^{-13}$). With $f(\xi)=\xi^{-3}\frac{d^3N_{\rm LO}}{d^3k}(\frac{k^+}{\xi},\frac{\mathbf k}{\xi})$ we have $f(1)=\frac{d^3N_{\rm LO}}{d^3k}(k^+,\mathbf k)$, and
\begin{equation}
\log\frac1{1-\xi_1}=\log\frac{k^++\Lambda}{\Lambda}=\log\frac{\vee}{\Lambda}-\log\frac{\vee}{k^+}+\log\Big(1+\frac{\Lambda}{k^+}\Big).
\end{equation}
So the pole \eqref{eq:pole} of Row IV is, exactly,
\begin{align}
\frac{d^3N_{\rm IV}}{d^3k}\bigg|_{\rm pole}={}&\frac{\alpha_s}{2\pi}\frac{d_\perp}{2}\frac{1}{\bar\epsilon}\bigg\{\int_{\xi_0}^1\frac{d\xi}{\xi^3}P^{(r)}_{gg,+}(\xi)\LO\Big(\frac{k^+}{\xi},\frac{\mathbf k}{\xi}\Big)\notag\\
&+2N_c\Big[\log\frac{\vee}{\Lambda}-\log\frac{\vee}{k^+}+\log\Big(1+\frac{\Lambda}{k^+}\Big)\Big]\LO(k^+,\mathbf k)-R_\Lambda\bigg\},
\label{eq:poleplus}
\end{align}
with
\begin{equation}
R_\Lambda\equiv\int_{\xi_1}^1d\xi\Big[P^{(r)}_{gg}(\xi)\,\xi^{-3}\LO\Big(\frac{k^+}{\xi},\frac{\mathbf k}{\xi}\Big)-\frac{2N_c}{1-\xi}\LO(k^+,\mathbf k)\Big]=-2N_c\frac{\Lambda}{k^+}\Big[f(1)+f'(1)\Big]+O\Big(\frac{\Lambda^2}{k^{+2}}\Big).
\end{equation}
$R_\Lambda$ is regular. The expansion was checked numerically: $R_\Lambda/\delta=$ $-1.97000$, $-1.97021$, $-1.97025$ for $\delta=10^{-2},10^{-3},10^{-4}$, vs.\ $-1.97025$. It is not of the form (number)$\times$LO.
\begin{itemize}
\item The first line of Eq.~\eqref{eq:poleplus} is the real part of $P_{gg}\otimes{\rm LO}$.
\item The second line is a $\delta(1-\xi)$ term: the soft logarithm $2N_c\log\frac{k^++\Lambda}{\Lambda}$. It cancels against the virtual rows of Group I (Sec.~\ref{sec:delta}).
\item The full Row IV, Eq.~\eqref{eq:conv}, is rewritten the same way with $f(\xi)=\xi^{-3}\mathcal L(\frac{k^+}{\xi},\frac{\mathbf k}{\xi})$. The $\mathcal L_D$ term is regular at $\xi=1$ and needs no plus prescription.
\end{itemize}

%==========================================================================================
\section{The \texorpdfstring{$\overline{\rm MS}$}{MS-bar} fragmentation function}\label{sec:FF}
\paragraph{Factorization.} A hadron (or any final state) $h$ with momentum $(k^+,\mathbf k)$ is produced from a gluon with $(k^+,\mathbf k)/\xi$:
\begin{equation}
\frac{d^3N_h}{d^3k}(k^+,\mathbf k)=\int_{\xi_0}^1\frac{d\xi}{\xi^3}\,D^{(0)}_{h/g}(\xi)\,\frac{d^3N_g}{d^3k}\Big(\frac{k^+}{\xi},\frac{\mathbf k}{\xi}\Big),
\label{eq:fact}
\end{equation}
where $D^{(0)}_{h/g}$ is the bare fragmentation function and $\frac{d^3N_g}{d^3k}=\LO+\frac{d^3N_{\rm NLO}}{d^3k}$ is the gluon spectrum computed in these notes.

\paragraph{The pole of the gluon spectrum.} Sec.~\ref{sec:delta} adds the poles of Row IV, of the virtual rows of Group I, and of the charge renormalization of LO. Up to terms of order $\frac{\Lambda}{k^+}$ and $\frac{k^+}{\vee}$, the sum is
\begin{equation}
\frac{d^3N_g}{d^3k}\bigg|_{\rm pole}=\frac{\alpha_s}{2\pi}\frac{1}{\bar\epsilon}\int_{\xi_0}^1\frac{d\xi}{\xi^3}\,P_{gg}(\xi)\,\LO\Big(\frac{k^+}{\xi},\frac{\mathbf k}{\xi}\Big),
\label{eq:totpole}
\end{equation}
\begin{equation}
P_{gg}(\xi)=2N_c\Big[\frac{\xi}{(1-\xi)_+}+\frac{1-\xi}{\xi}+\xi(1-\xi)\Big]+\frac{11N_c}{6}\,\delta(1-\xi)\qquad(n_f=0).
\label{eq:Pgg}
\end{equation}
This is the full LO gluon splitting function. The coefficient $\frac{11N_c}{6}=\frac{\beta_0}{2}$ of $\delta(1-\xi)$ is also fixed independently by momentum conservation: $\int_0^1d\xi\,\xi\,P^{(r)}_{gg,+}=$ $-\frac{11N_c}{6}$ (\texttt{sympy}), so $\int_0^1d\xi\,\xi P_{gg}=0$ requires $\frac{11N_c}{6}\delta(1-\xi)$.

\paragraph{The counterterm.} Define
\begin{equation}
\boxed{\Gamma_{gg}(\xi)\equiv\delta(1-\xi)-\frac{\alpha_s}{2\pi}\frac{1}{\bar\epsilon}P_{gg}(\xi)=\delta(1-\xi)+\frac{\alpha_s}{2\pi}\frac{2}{\epsilon}P_{gg}(\xi)}
\label{eq:Gamma}
\end{equation}
(the second form is in your $\epsilon$, $d_\perp=2-\epsilon$; $\overline{\rm MS}$ subtracts only $\frac{2}{\epsilon}$, since $\mu_{\overline{MS}}$ already contains $\frac{e^{\gamma_E}}{4\pi}$). The renormalized fragmentation function $D_{h/g}(\xi,\mu)$ is defined by
\begin{equation}
D^{(0)}_{h/g}(\xi)=\int_\xi^1\frac{dz}{z}\,\Gamma_{gg}(z)\,D_{h/g}\Big(\frac{\xi}{z},\mu\Big)=D_{h/g}(\xi,\mu)-\frac{\alpha_s}{2\pi}\frac{1}{\bar\epsilon}\int_\xi^1\frac{dz}{z}P_{gg}(z)\,D_{h/g}\Big(\frac{\xi}{z},\mu\Big)+O(\alpha_s^2).
\label{eq:ren}
\end{equation}
Insert Eq.~\eqref{eq:ren} into Eq.~\eqref{eq:fact} and use the Jacobian identity (with $\xi=\zeta z$)
\begin{equation}
\int_{\xi_0}^1\frac{d\xi}{\xi^3}\Big[\int_\xi^1\frac{dz}{z}A(z)D\Big(\frac\xi z\Big)\Big]F\Big(\frac{k}{\xi}\Big)=\int_{\xi_0}^1\frac{d\zeta}{\zeta^3}D(\zeta)\int_{\xi_0/\zeta}^1\frac{dz}{z^3}A(z)\,F\Big(\frac{k}{\zeta z}\Big)
\label{eq:jac}
\end{equation}
(numerically, rel.\ $4\times10^{-14}$). Then
\begin{equation}
\frac{d^3N_h}{d^3k}=\int_{\xi_0}^1\frac{d\xi}{\xi^3}\,D_{h/g}(\xi,\mu)\,\frac{d^3N^{\overline{MS}}_g}{d^3k}\Big(\frac{k^+}{\xi},\frac{\mathbf k}{\xi};\mu\Big),\qquad
\frac{d^3N^{\overline{MS}}_g}{d^3k}(k;\mu)\equiv\int\frac{dz}{z^3}\,\Gamma_{gg}(z)\,\frac{d^3N_g}{d^3k}\Big(\frac{k}{z}\Big).
\end{equation}
The counterterm $-\frac{\alpha_s}{2\pi}\frac{1}{\bar\epsilon}P_{gg}\otimes{\rm LO}$ cancels Eq.~\eqref{eq:totpole}, so $\frac{d^3N^{\overline{MS}}_g}{d^3k}$ is finite.

\paragraph{Row IV after the subtraction ($d_\perp=2-\epsilon$).} With Eq.~\eqref{eq:CDR}, the part of Row IV that stays in $\frac{d^3N^{\overline{MS}}_g}{d^3k}$ is Eq.~\eqref{eq:conv} without its $\frac{1}{\bar\epsilon}$:
\begin{equation}
\frac{d^3N^{\overline{MS}}_{\rm IV}}{d^3k}=\frac{\alpha_s}{2\pi}\int_{\xi_0}^{\xi_1}\frac{d\xi}{\xi^3}\Big\{-P^{(r)}_{gg}(\xi)\,\frac{d^3N_\ell}{d^3k}\Big(\frac{k^+}{\xi},\frac{\mathbf k}{\xi}\Big)+2N_c\,\xi(1-\xi)\Big[\LO-2\frac{d^3N_D}{d^3k}\Big]\Big(\frac{k^+}{\xi},\frac{\mathbf k}{\xi}\Big)\Big\}.
\end{equation}
The logarithm $\ell=\log\frac{(y'-y)^2\mu^2}{c_0^2}$ vanishes at $\mu=c_0/|y'-y|$, the natural scale of the fragmentation logarithm. With $|y'-y|\sim1/|\mathbf k|$, this is $\mu\sim|\mathbf k|$.

%==========================================================================================
\section{DGLAP}\label{sec:dglap}
\paragraph{From the $\mu$ independence of the cross section.} $\frac{d^3N_h}{d^3k}$ does not depend on $\mu$.
\begin{itemize}
\item In $\frac{d^3N^{\overline{MS}}_g}{d^3k}$ every pole came with its logarithm, $\bar E=\frac1{\bar\epsilon}-\log\frac{L_\perp^2\mu^2}{c_0^2}$, with $L_\perp=y'-y$ in Row IV, $x-w$ in Row II and $y-z$ in Row V 2b. After the subtraction, $\frac{\partial}{\partial\log\mu^2}$ gives $-1$ times the residue. Summed over the rows, the residues are $P^{(r)}_{gg,+}$ plus the $\delta(1-\xi)$ terms of Sec.~\ref{sec:delta} without the charge renormalization. The soft logarithms and the $\frac{11}{3}$ cancel there.
\item The LO term carries $\alpha_s(\mu)$, with $\frac{d\alpha_s}{d\log\mu^2}=-\frac{\beta_0}{4\pi}\alpha_s^2$. This gives $-\frac{\alpha_s}{2\pi}\frac{\beta_0}{2}{\rm LO}$.
\end{itemize}
Together,
\begin{equation}
\frac{\partial}{\partial\log\mu^2}\frac{d^3N^{\overline{MS}}_g}{d^3k}(k^+,\mathbf k;\mu)=-\frac{\alpha_s}{2\pi}\int_{\xi_0}^1\frac{dz}{z^3}P_{gg}(z)\,\LO\Big(\frac{k^+}{z},\frac{\mathbf k}{z}\Big)+O(\alpha_s^2).
\end{equation}
By Eq.~\eqref{eq:jac}, $\mu$ independence of $\frac{d^3N_h}{d^3k}$ requires
\begin{equation}
\boxed{\mu^2\frac{\partial}{\partial\mu^2}D_{h/g}(\xi,\mu)=\frac{\alpha_s(\mu)}{2\pi}\int_\xi^1\frac{dz}{z}\,P_{gg}(z)\,D_{h/g}\Big(\frac{\xi}{z},\mu\Big)}
\label{eq:DGLAP}
\end{equation}
with $P_{gg}$ of Eq.~\eqref{eq:Pgg}. This is the time-like DGLAP equation for gluon fragmentation (pure glue, $n_f=0$).

\paragraph{From the $\mu$ independence of $D^{(0)}$.} The same follows from Eq.~\eqref{eq:ren}. In $d_\perp=2-\epsilon$, $g^2$ has mass dimension $\epsilon$, so $\frac{d\alpha_s}{d\log\mu^2}=-\frac{\epsilon}{2}\alpha_s+O(\alpha_s^2)$. Then
\begin{equation}
\frac{\partial}{\partial\log\mu^2}\Big[\frac{\alpha_s}{2\pi}\frac{1}{\bar\epsilon}\Big]=\frac{\alpha_s}{2\pi}\Big(-\frac\epsilon2\Big)\Big(-\frac2\epsilon\Big)=\frac{\alpha_s}{2\pi}
\end{equation}
(\texttt{sympy}: ratio $1$), and $\frac{dD^{(0)}}{d\log\mu^2}=0$ gives Eq.~\eqref{eq:DGLAP}.

\paragraph{Solution to first order.}
\begin{equation}
D_{h/g}(\xi,\mu)=D_{h/g}(\xi,\mu_0)+\frac{\alpha_s}{2\pi}\log\frac{\mu^2}{\mu_0^2}\int_\xi^1\frac{dz}{z}P_{gg}(z)\,D_{h/g}\Big(\frac{\xi}{z},\mu_0\Big)+O(\alpha_s^2).
\end{equation}
With $\mu\sim|\mathbf k|$, the logarithms $\ell$ of Row IV are small, and the large collinear logarithms are resummed by Eq.~\eqref{eq:DGLAP}.

%==========================================================================================
\section{The \texorpdfstring{$\delta(1-\xi)$}{delta(1-xi)} part: Row IV and Group I}\label{sec:delta}
Row IV alone has no $\frac{11N_c}{6}\delta(1-\xi)$. Its $\delta(1-\xi)$ term is the soft logarithm of Eq.~\eqref{eq:poleplus}. The other poles proportional to LO are in Group I, Row II and Row V (part 2b, UV). In \texttt{Group\_I\_corrected.pdf}, with the exact brackets (G5a, G5d), they are:
\begin{itemize}
\item \textbf{Row II} (per ordering):
\begin{equation*}
\frac{1}{(2\pi)^3}\frac{g^4N_c}{32\pi^4}\frac1{k^+}\int_{w,w',x,x'}e^{-ik\cdot(w'-w)}\frac{(x'-w')\cdot(x-w)}{(x'-w')^2(x-w)^2}\,\bar E_{x-w}\,B_{\rm II}\,[U(x')-U(w')][U(w)-U(x)]\rho\rho,
\end{equation*}
\begin{equation*}
B_{\rm II}=\frac{11}{3}+\log\frac{\Lambda^2}{k^{+2}-\Lambda^2}+\log\Big(1-\frac{k^+}{\vee}\Big)-8\delta+2\delta^2-\frac43\delta^3 .
\end{equation*}
The colour factor is $-\OO$ with $y\to w$, $y'\to w'$ (numerically: ratio $-1$). The prefactor relative to LO is $\frac{\alpha_s}{2\pi}\cdot$$\frac{N_c}{2}$. So, per ordering, Row II $\to-\frac{\alpha_s}{2\pi}\frac{N_c}{2}\bar E\,B_{\rm II}\times$LO.
\item \textbf{Row V, 2b UV} (per ordering):
\begin{equation*}
\frac{1}{(2\pi)^3}\frac{g^4N_c}{16\pi^4}\frac1{k^+}\,B_{\rm V}\int_{y,z,x',w'}e^{-ik\cdot(w'-z)}\,E_{y-z}\,\frac{(x'-w')\cdot(y-z)}{(x'-w')^2(y-z)^2}[U(x')-U(w')][U(y)-U(z)]\rho\rho,
\end{equation*}
\begin{equation*}
B_{\rm V}=2\log\frac{\vee}{\Lambda}-2\log\frac{\vee}{k^+}+2\log\Big(1-\frac{\Lambda}{k^+}\Big)-\frac{11}{6}+4\delta-\delta^2+\frac23\delta^3 .
\end{equation*}
The colour factor is $+\OO$ with $x\to y$, $y\to z$, $y'\to w'$ (ratio $+1$). The prefactor is $\frac{\alpha_s}{2\pi}\cdot$$N_c$. With $E=-\bar E$, per ordering Row V 2b $\to-\frac{\alpha_s}{2\pi}N_c\bar E\,B_{\rm V}\times$LO.
\item \textbf{Complex conjugate.} Group I is $\mathcal C\cdot\mathcal V+{\rm c.c.}$ (Eq.~(CS1)), so both rows count twice. Row IV is the square $4\mathfrak C'^\dagger\mathfrak C'$ and counts once.
\item \textbf{Charge renormalization.} With $\alpha_0=\alpha_s(\mu)\big[1-\frac{\alpha_s\beta_0}{4\pi}\frac1{\epsilon_{\rm std}}\big]=\alpha_s(\mu)\big[1+\frac{\alpha_s\beta_0}{4\pi}\frac1{\bar\epsilon}\big]$, ${\rm LO}(\alpha_0)={\rm LO}(\alpha_s(\mu))+\frac{\alpha_s}{2\pi}\frac{\beta_0}{2}\frac{1}{\bar\epsilon}{\rm LO}$.
\end{itemize}
\begin{table}[H]\centering\small
\renewcommand{\arraystretch}{1.3}
\begin{tabular}{@{}lll@{}}\toprule
source & weight & coefficient of $\frac{\alpha_s}{2\pi}\frac{1}{\bar\epsilon}N_c\,{\rm LO}(k^+,\mathbf k)$\\\midrule
Row IV, Eq.~\eqref{eq:poleplus} & 1 & $2\log\frac{1+\delta}{\delta}\ \big(=2\delta Y-2\log\frac{\vee}{k^+}+2\log(1+\delta)\big)$; and $-R_\Lambda$ (not $\propto$ LO)\\
Row II & $\times2$ & $-B_{\rm II}=-\frac{11}{3}+2\log\frac1\delta+\log(1-\delta^2)-\log(1-\xi_0)+8\delta-2\delta^2+\frac43\delta^3$\\
Row V, 2b UV & $\times2$ & $-2B_{\rm V}=\frac{11}{3}-4\log\frac1\delta-4\log(1-\delta)-8\delta+2\delta^2-\frac43\delta^3$\\
charge renormalization & & $+\frac{11}{6}$\\\midrule
sum & & $\frac{11}{6}+3\log\frac{1+\delta}{1-\delta}-\log(1-\xi_0)$; and $-R_\Lambda$\\\bottomrule
\end{tabular}
\end{table}
(\texttt{sympy}: sum minus the last line $=$ $0$; here $\log\frac1\delta=\delta Y-\log\frac{\vee}{k^+}$.)
\begin{itemize}
\item \textbf{The soft logarithms cancel:} $\delta Y$ appears with $+2$ (Row IV), $+2$ (Row II) and $-4$ (Row V 2b). This holds only with the weights $1:2:2$. So it also checks the relative normalization of Row IV and of the c.c.\ in Group I.
\item \textbf{The $\frac{11}{3}$ cancel} between Row II and Row V 2b (per ordering, as found in \texttt{Part2b\_MSbar.pdf}), and so do the powers $8\delta-2\delta^2+\frac43\delta^3$.
\item \textbf{The $\frac{\beta_0}{2}\delta(1-\xi)$ of $P_{gg}$ is the charge renormalization of LO.} This resolves the open point of \texttt{Part2b\_MSbar.pdf} (and of the subsection ``Renormalization-group bookkeeping'' in \texttt{sec8\_rc\_dglap.tex} of the paper draft): no further row is needed to supply $\frac{\alpha_s\beta_0}{4\pi}\frac{1}{\epsilon_{\rm std}}$. With one $\epsilon$ for all poles, the pole of the coupling counterterm is not cancelled; it becomes the $\delta(1-\xi)$ term of the fragmentation counterterm \eqref{eq:Gamma}. Together with Row IV it gives the complete $P_{gg}$, Eq.~\eqref{eq:totpole}.
\item \textbf{Exact remainder.} For finite $\Lambda/k^+$ and $k^+/\vee$, the sum leaves
\begin{equation}
\frac{\alpha_s}{2\pi}\frac{1}{\bar\epsilon}\Big\{N_c\Big[3\log\frac{1+\delta}{1-\delta}-\log\Big(1-\frac{k^+}{\vee}\Big)\Big]\LO(k^+,\mathbf k)-R_\Lambda\Big\},\qquad 3\log\frac{1+\delta}{1-\delta}=6\delta+2\delta^3+\dots
\end{equation}
This is of order $\frac{\Lambda}{k^+}$ and $\frac{k^+}{\vee}$. It is not of the form $P\otimes{\rm LO}$, since $R_\Lambda$ involves $f'(1)$, so it cannot go into $D_{h/g}$.
\begin{itemize}
\item It comes from the different longitudinal limits of the rows: $p^+>\Lambda$ in Row IV; $\Lambda<p^+<k^+-\Lambda$ in Row V 2b; the loop up to $\vee$ in Row II.
\item Either it cancels against rows not included here, or the limits must be made consistent.
\item This is still open, as was already noted in \texttt{Part2b\_MSbar.pdf} for the poles multiplying $\log(1-\frac{k^+}{\vee})$ and the powers of $\frac{\Lambda}{k^+}$.
\end{itemize}
\end{itemize}

%==========================================================================================
\section{Checks}
Script \texttt{frag\_checks.py}, results in \texttt{frag\_checks.json}.
\begin{center}\small
\begin{tabular}{@{}p{0.62\textwidth}p{0.32\textwidth}@{}}\toprule
check & result\\\midrule
your Eq.~\eqref{eq:user} minus Eq.~(1.5) of \texttt{RowIV\_dperp.pdf} (\texttt{sympy}) & $0$\\
conversion factor $\frac{\alpha_s}{2\pi}\frac{N_c}{\xi^3}$ (\texttt{sympy}) & $0$\\
convolution form \eqref{eq:conv} minus compact form (\texttt{sympy}) & $0$\\
$d_\perp=2-\epsilon$: $\frac{d_\perp\bar E-2}{2}-\big(\frac1{\bar\epsilon}-\ell\big)$ at $O(\epsilon^0)$ (\texttt{sympy}) & $0$\\
plus prescription with exact limits \eqref{eq:plusid}, 24 random cases & rel.\ $1\times10^{-13}$\\
$R_\Lambda/\delta\to-2N_c[f(1)+f'(1)]$ & $-1.97000$, $-1.97021$, $-1.97025$ vs.\ $-1.97025$\\
momentum sum rule, $\delta(1-\xi)$ coefficient (\texttt{sympy}) & $\frac{11N_c}{6}\delta(1-\xi)$\\
Jacobian identity \eqref{eq:jac} & rel.\ $4\times10^{-14}$\\
$\frac{\partial}{\partial\log\mu^2}\big[\frac{\alpha_s}{2\pi}\frac{1}{\bar\epsilon}\big]\big/\frac{\alpha_s}{2\pi}$ (\texttt{sympy}) & $1$\\
colour factors of Row II and Row V 2b relative to $\OO$ & $-1$, $+1$\\
$\delta(1-\xi)$ tally, exact in $\delta$ and $\xi_0$ (\texttt{sympy}) & $0$\\\bottomrule
\end{tabular}
\end{center}
\end{document}
